\documentclass[10pt,twocolumn,letterpaper]{article}
\usepackage[pagenumbers]{cvpr}
\usepackage{graphicx}
\usepackage{amsmath,amssymb}
\usepackage{booktabs}
\usepackage{xspace}
\definecolor{cvprblue}{rgb}{0.21,0.49,0.74}
\usepackage[breaklinks,colorlinks,allcolors=cvprblue]{hyperref}
\newcommand{\ours}{G$\times$S\xspace}
\newcommand{\TODO}[1]{\textbf{\color{red}[#1]}}
\def\confName{CVPR}
\def\confYear{2026}

\title{\ours: Decoupling \emph{How} to Grasp from \emph{What} to Grasp\\for Language-Driven Grasp Detection}
\author{Erhan Lai\\{\small\url{https://github.com/RyleHan/gxs-grasp}}}

\begin{document}
\maketitle

\begin{abstract}
Language-driven grasp detection asks for a grasp rectangle that satisfies a natural-language instruction.
Current methods inject the instruction into the whole grasp network, so the text also has to shape
the grasp geometry, which it does not determine.
We factorise the grasp quality map into a language-agnostic \emph{graspability} map $G$ and an
instruction-conditioned \emph{selection} map $S$.
$G$ is a GR-ConvNet trained on the union of all grasps annotated in an image;
$S$ compares frozen dense CLIP features with the instruction and only has to choose among graspable
regions. Both are supervised with labels that Grasp-Anything++ already provides: every image carries
several instructions, and their grasps act as each other's negatives, so no masks are needed.
On the official test split \ours raises the success rate from \TODO{x} to \TODO{y}, and the rate at
which both grasps of an image pair are correct from \TODO{a} to \TODO{b}.
\end{abstract}

\section{Introduction}
A robot asked to \emph{``grip the keychain by its ring''} has to answer two questions:
which part of the image the sentence refers to, and how a parallel gripper should close on it.
Grasp-Anything++~\cite{vuong2024lgd} made this task learnable at scale, pairing about a million
synthetic scenes with grasp instructions. Methods trained on it~\cite{vuong2024lgd,vo2024maskgrasp}
condition the entire detector on the text.

Two properties of the data suggest a simpler structure.
First, the opening width and angle of a grasp at a given pixel follow from the object's shape;
the sentence only decides \emph{where} to grasp.
Second, instructions come in groups: one image carries a grasp instruction for each annotated object,
and these point to different regions. The union of their grasps says where the image can be grasped
at all; their differences say what each sentence adds.

We build on both observations. \ours predicts $Q = G\cdot S$, where $G$ ignores the language and $S$
selects among the graspable regions. Our contributions are:
(i) a factorised detector whose geometry heads never see the text;
(ii) a sibling-contrast loss that supervises selection with the other instructions of the same image,
without any masks; and
(iii) a paired metric that a model ignoring the instruction cannot pass, together with an analysis of
the data showing that most part-level labels do not change with the part named in the instruction.

\section{Related Work}
\textbf{Language-driven grasping.}
CLIPort~\cite{shridhar2021cliport} fuses CLIP~\cite{radford2021clip} features into dense
pick-and-place affordances. CROG~\cite{tziafas2023crog} extends CLIP-based referring segmentation~\cite{wang2022cris}
with grasp maps. On Grasp-Anything++~\cite{vuong2024graspanything,vuong2024lgd}, LGD casts grasp detection as
conditional diffusion, and mask-guided methods~\cite{vo2024maskgrasp,bhat2025maplegrasp} use object
masks to focus the grasp features. Unlike these, we multiply at the \emph{output} rather than the
feature level, keep geometry language-free, and need no mask annotations.

\textbf{Dense CLIP features.}
CLIP patch tokens are poorly localised; MaskCLIP~\cite{zhou2022maskclip} uses the value embedding of the
last attention layer and ClearCLIP~\cite{lan2024clearclip} drops its residual and feed-forward paths.
We adopt this combination without training.

\section{Method}
Given an RGB image $I$ and an instruction $\ell$, we predict $g=(x,y,w,h,\theta)$.

\textbf{Graspability $G$.}
A GR-ConvNet~\cite{kumra2020grconvnet} maps $I$ to five $224{\times}224$ maps: quality logits $q$,
$\cos2\theta$, $\sin2\theta$ (a parallel gripper is symmetric under $180^\circ$), opening $w$ and jaw
size $h$. None of them receive $\ell$. For an image with instructions $\ell_1,\dots,\ell_K$ and grasp
sets $\mathcal{G}_1,\dots,\mathcal{G}_K$, targets are rendered from $\bigcup_k\mathcal{G}_k$:
\begin{equation}
\mathcal{L}_G=\mathrm{BCE}(q,\hat q)+\sum_{m\in\{\cos,\sin,w,h\}}\big\lVert (m-\hat m)\odot\hat q\big\rVert_{\text{sL1}},
\end{equation}
where $\hat q$ marks the centre third of each rectangle. Supervising geometry only where a grasp
exists matters: with GR-ConvNet's dense regression the predicted opening collapsed to
$\sim$15\,px ($\sim$80\,px ground truth) even when overfitting 70 scenes.

\textbf{Selection $S$.}
Frozen CLIP ViT-B/16 at $448{\times}448$ gives $28{\times}28$ patch features $v_p$ (value-only last
block); $t$ is the CLIP text embedding. With zero-initialised residual adapters $\phi,\psi$,
\begin{equation}
S(p\mid\ell)=\sigma\!\big(\tau\,\langle\phi(v_p),\psi(t)\rangle+b\big).
\end{equation}
For instruction $\ell_k$, patches covered by $\mathcal{G}_k$ are positives, patches covered by a
sibling $\mathcal{G}_{j\neq k}$ are negatives, and the rest are ignored, since background is already
suppressed by $G$. We use a class-balanced BCE over these patches.

\textbf{Inference.}
$Q=\sigma(q)\cdot S\!\uparrow$ ($S$ bilinearly upsampled); the grasp centre is $\arg\max Q$ and
$\theta,w,h$ are read from $G$'s maps at that pixel. The backbone runs once per image;
a new instruction costs one dot product.

\section{Experiments}
\textbf{Data.} We keep the official LGD test splits (seen / unseen categories).
The official training split contains only \TODO{n} multi-object scenes, so we add \TODO{n}
multi-object scenes from the full release, requiring that (R1) no scene appears in a test split,
(R2) every object name belongs to the official training vocabulary, and (R3) no two objects share a
name. In the resulting training set, \TODO{r}\% of part pairs of the same object have identical grasp labels,
so language mostly has to pick the object.

\textbf{Metrics.} \emph{Success}: IoU\,$>0.25$ and $|\Delta\theta|<30^\circ$ with a ground-truth grasp
of the instruction~\cite{vuong2024lgd}. \emph{Paired}: for each pair of objects in a multi-object
test image, both grasps must succeed; ignoring the instruction gives almost zero.
\emph{Select}: the predicted centre lies on the target object.

\textbf{Results.}
\TODO{table + discussion}

\section{Conclusion}
Separating how to grasp from what to grasp lets each part learn from the supervision it actually
needs. The selection map inherits CLIP's open vocabulary, while geometry is learned once per image.
The main limitation is CLIP's coarse $28{\times}28$ grid, which limits part-level selection;
the labels of Grasp-Anything++ rarely distinguish parts anyway.

{\small
\bibliographystyle{ieeenat_fullname}
\bibliography{main}
}
\end{document}
